% SPDX-FileCopyrightText: Copyright (c) 2016-2026 Objectionary.com
% SPDX-License-Identifier: MIT

\section{Reading the Specification}\label{sec:preamble}

This section is a brief guide to the rest of the paper.
It states what each chapter and appendix contains, suggests reading
  orders for three audiences, distinguishes normative from informative
  material, and fixes the typographic and mathematical conventions used
  throughout.

\subsection{What Each Chapter Contains}\label{sec:preamble:chapters}

\begin{description}
  \item[\cref{sec:overview}, \emph{Informal Overview}.]
    A tour of EO's constructs paired with their \(\varphi\)-calculus
    translations. Informal: examples drive the prose, no construct is
    given its full surface contract here. Read this first.
  \item[\cref{sec:lexical}, \emph{Lexical Layer}.]
    Encoding, indentation, token catalogue, numeric and string literal
    grammars, the multi-line bytes form. Normative.
  \item[\cref{sec:syntax}, \emph{Syntax}.]
    Per-line shapes (\cref{sec:syntax:lines}), name suffixes
    (\cref{sec:syntax:names}), multi-line constructs
    (\cref{sec:syntax:multi-line}), the XMIR output mapping
    (\cref{sec:syntax:xmir}), and the twelve outer kinds
    (\cref{sec:syntax:outer-kinds}). Normative.
    The line-shape decision table at the head of \cref{sec:syntax:lines}
    is the entry point for classifying any source line.
  \item[\cref{sec:wellformedness}, \emph{Well-formedness Rules}.]
    The cross-line rules that decide which syntactically-recognisable
    programs are legal: notation consistency
    (\cref{sec:wf:notation}), naming inside formations
    (\cref{sec:wf:naming}), atoms and tests (\cref{sec:wf:atoms}),
    comment attachment (\cref{sec:wf:comments}), blank-line policy
    (\cref{sec:wf:blanks}), inline bindings (\cref{sec:wf:bindings}).
    Normative.
  \item[\cref{sec:semantics}, \emph{Semantics via \(\varphi\)-Calculus}.]
    The translation function \(\llbracket\,\cdot\,\rrbracket\) from
    well-formed EO expressions to \(\varphi\)-calculus expressions,
    construct-by-construct.
    Reserved attributes, assets (\(\Delta\), \(\lambda\)),
    dataization, meta directives, and aliases also live here.
    Normative; reads against the \(\varphi\)-calculus surface defined
    by \citet{bugayenko2021eolang}.
  \item[\cref{sec:related}, \emph{Related Work}.]
    EO's position among pure-OO languages, mainstream OO languages, and
    object calculi. Informative.
  \item[\cref{sec:future}, \emph{Future Studies}.]
    Directions that build on this specification. Informative.
  \item[\cref{app:outer-kinds}.]
    At-a-glance reference table for the twelve outer kinds. Normative.
  \item[\cref{app:examples}.]
    Worked-through EO programs covering most surface constructs in
    combination, each paired with prose and selected \(\varphi\)-form
    translations. Informative.
\end{description}

\subsection{Reading Orders}\label{sec:preamble:orders}

\begin{description}
  \item[Curious reader.]
    \Cref{sec:overview} is self-contained; finish there and read the
    rest only if the syntax appears arbitrary.
  \item[Practitioner writing EO programs.]
    \Cref{sec:overview} for orientation, \cref{sec:syntax} for the
    surface forms, \cref{sec:wellformedness} for the rules that decide
    which compositions are accepted, \cref{app:examples} for canned
    patterns.
  \item[Compiler or tool implementer.]
    Read in document order: \cref{sec:lexical} drives lexing,
    \cref{sec:syntax} drives the line classifier and outer-kind
    assignment, \cref{sec:wellformedness} drives the cross-line
    validator, \cref{sec:semantics} fixes the translation contract.
    \Cref{app:outer-kinds} is the cross-reference table to keep open
    on the side.
  \item[Theorist.]
    \Cref{sec:semantics} is the bridge to the \(\varphi\)-calculus;
    read it alongside \citet{bugayenko2021eolang}.
\end{description}

\subsection{Normative versus Informative}\label{sec:preamble:status}

A chapter or appendix is \emph{normative} if a conforming implementation
  is obliged to honour every rule it states.
A chapter is \emph{informative} if it expounds, motivates, or
  illustrates but does not add obligations.
\Cref{sec:lexical,sec:syntax,sec:wellformedness,sec:semantics} and
  \cref{app:outer-kinds} are normative.
The abstract, \cref{sec:introduction,sec:overview,sec:related,sec:future},
  and \cref{app:examples} are informative.
Where a normative chapter introduces an example, the example is
  illustrative; the prose around it states the normative rule.

\subsection{Typographic and Notational Conventions}\label{sec:preamble:conventions}

\paragraph{Code and tokens.}
EO source and lexical tokens are typeset in a fixed-width face:
  \ff{[a b] > foo}, \ff{NAME}, \ff{BYTES}.
The same face is used for canonical objects of the standard library
  (\ff{Q.io.stdout}, \ff{Q.tt.sprintf}) and for XMIR markup
  (\ff{<o>}, \ff{@name}).

\paragraph{\texorpdfstring{\(\varphi\)-calculus}{phi-calculus} terms.}
The translation function is \(\llbracket\,\cdot\,\rrbracket\).
Greek-letter names from the \(\varphi\)-calculus
  (\(\Phi\), \(\xi\), \(\varphi\), \(\rho\), \(\Delta\), \(\lambda\))
  are typeset in math italic, with \(\Phi\) reserved for the global
  root namespace, \(\xi\) for the current scope, \(\varphi\) for the
  decoratee, \(\rho\) for the parent, and \(\Delta\), \(\lambda\) for
  the two reserved asset slots.
Positional attribute labels are \(\alpha_0, \alpha_1, \ldots\),
  generated by the translation in source order.

\paragraph{Cross-references.}
Internal references use the \texttt{cleveref} \texttt{\textbackslash cref}
  macro: \enquote{see \cref{sec:wf:atoms}} rather than
  \enquote{see Section X.Y}.
A cross-reference to an entire chapter cites its top-level label
  (\cref{sec:syntax}); references to subsections cite the more specific
  label (\cref{sec:syntax:outer-kinds}).

\paragraph{Examples.}
Examples are presented as fenced \ff{ffcode} blocks.
Line comments inside an example use \ff{//} and are not part of EO
  source; they are typographical annotations of the example.
Where an example is illegal, the offending line is annotated
  \enquote{\ff{rejected}} in its trailing comment.

\paragraph{Indent depths.}
Indent depths are stated in spaces (\enquote{indent 4}) when the
  absolute count matters, or as levels (\enquote{indent level 1},
  meaning indent 2) when only the nesting depth matters.

\paragraph{Outer kinds and openness.}
The vocabulary of \emph{outer kinds} and \emph{openness} from
  \cref{sec:syntax:outer-kinds} is used throughout
  \cref{sec:wellformedness,sec:semantics} and the appendices;
  the reference table is \cref{app:outer-kinds}.
